\section{The support-size boundary}
\label{sec:boundary}

\begin{proposition}\label{b-small-support}
For every nonempty integer set $F$ with $|F|\le3$ and every nonnegative
kernel $k$ supported on $F$, one has $N_F(k)^4\le T_F(k)$.
\end{proposition}

\begin{proof}
Sets with at most two points are Sidon, as is any three-point set that is
not an arithmetic progression.  Proposition~\ref{c-sidon} applies to all
of them.  For a three-term progression, complete it to a four-term
progression and extend both $x$ and $k$ by zero.  This preserves ordinary
convolution, the weighted objective, and $T_F(k)$.  The four-term-progression
inequality proved above therefore gives the desired bound for every feasible
$x$, and then for $N_F(k)$.
\end{proof}

For a singleton $F=\{p\}$ the bound is equality:
$N_F(k)=k(p)$ and $T_F(k)=k(p)^4$.  The corresponding max-convolution
inequality is also equality for every nonzero $f,g$ allowed in
Corollary~\ref{cor:h4}, since
\[
 \|f\star g\|_1=g(p)\|f\|_1,\qquad
 \|f\star\widetilde{gk}\|_1=g(p)k(p)\|f\|_1.
\]
Thus the strict four-point statement does not extend to singleton supports.

The preceding proposition, the classification in
Proposition~\ref{c-classification}, and the five four-point cases prove
the positive part of Theorem~\ref{thm:main}.  The following example proves
its support-size assertion.

\begin{proposition}\label{b-five-point}
Let $F=\{0,2,7,8,11\}$ and $k=1_F$.  Then
\[
 N_F(k)=\frac52,\qquad T_F(k)=39,\qquad
 N_F(k)^4-T_F(k)=\frac1{16}.
\]
\end{proposition}

\begin{proof}
The fifteen unordered pair sums, including diagonal pairs, are
\[
 0,2,4,7,8,9,10,11,13,14,15,16,18,19,22.
\]
They are distinct, so $F$ is Sidon.  For $x=1_F/\sqrt2$, ordinary
convolution equals $1/2$ at diagonal sums and $1$ at off-diagonal sums.
Consequently $x*x\le1$ and $\sum_{p\in F}x(p)^2=5/2$.
Equation~\eqref{c-sidon-norm}, which holds for Sidon sets of any size,
also supplies the upper bound:
$N_F(1_F)=\max\{5/2,(5+3)/4\}=5/2$.

To count $4F$, put $E=\{0,2,7,8\}$ and write
$[r,s]_{\mathbb Z}=\{r,r+1,\ldots,s\}$ for integer endpoints $r\le s$.
For $h\ge0$, partition the summands in $hE$ by the number $m$ drawn
from $\{7,8\}$.  Their sum ranges from $7m$ to $8m$ in unit steps,
while the remaining summands from $\{0,2\}$ contribute $2j$ for
$0\le j\le h-m$.  Hence
\begin{equation}\label{b-small-sumsets}
 hE=\bigcup_{m=0}^{h}\ \bigcup_{j=0}^{h-m}
           [7m+2j,8m+2j]_{\mathbb Z}.
\end{equation}
In particular,
\[
 4E=\{0,2,4\}\cup[6,26]_{\mathbb Z}\cup[28,32]_{\mathbb Z}.
\]
Now partition $4F$ by the number $\ell$ of summands equal to $11$.
Formula~\eqref{b-small-sumsets} gives the new elements at each step:
\begin{center}
\begin{tabular}{clc}
\toprule
Layer & Elements not in preceding layers & Number \\
\midrule
$4E$ & $\{0,2,4\}\cup[6,26]_{\mathbb Z}\cup[28,32]_{\mathbb Z}$ & $29$ \\
$11+3E$ & $\{27,33,34,35\}$ & $4$ \\
$22+2E$ & $\{36,37,38\}$ & $3$ \\
$33+E$ & $\{40,41\}$ & $2$ \\
$\{44\}$ & $\{44\}$ & $1$ \\
\bottomrule
\end{tabular}
\end{center}
Thus
\[
 4F=\{0,2,4\}\cup[6,38]_{\mathbb Z}\cup\{40,41,44\},
 \qquad |4F|=39.
\]
The fourfold max-convolution of an indicator is the indicator of its
fourfold sumset, so $T_F(1_F)=39$.  Finally,
$(5/2)^4=625/16=39+1/16$.
\end{proof}

The example shows failure of the sufficient covering bound
$N_F(k)^4\le T_F(k)$ with both $x$ and $k$ positive throughout a
five-point set.  A counterexample to the reflected max-convolution
inequality would additionally require functions $f,g$ for which the ratio
in \eqref{eq:ratio} exceeds $T_F(k)^{1/4}$; this construction supplies
no such functions.

\section{Sharp unweighted covering constants}
\label{sec:constants}

The five pair-sum patterns have different unweighted norms.  Their values,
with one maximizing squared profile in the natural coordinate order, are
listed below.
\begin{center}
\begin{tabular}{lcl}
\toprule
Four-point pattern & $N_F(1_F)$ & Maximizing squared profile \\
\midrule
Sidon & $2$ & $(1/2,1/2,1/2,1/2)$ \\
Isolated progression, $(0,a,2a,b)$ & $7/4$ & $(1/4,1/4,1/4,1)$ \\
Two progressions, $(0,1,2,4)$ & $105/64$ & $(9/64,1/4,1/4,1)$ \\
Parallelogram, $(0,a,b,a+b)$ & $25/16$ & $(1,1/4,1/4,1/16)$ \\
Four-term progression & $3/2$ & $(1,1/4,0,1/4)$ \\
\bottomrule
\end{tabular}
\end{center}

For Sidon supports, \eqref{c-sidon-norm} gives
$\max\{2,7/4\}=2$.  The isolated-progression profile sums are
$105/64$ twice, $3/2$ four times, and $7/4$ twice.  The sixteen
parallelogram profile sums are $3/2$ for the twelve profiles supported
on three corners and $25/16$ for the four tensor profiles.  For a
four-term progression, the five profile sums before reversal are
$3/2,3/2,381/256,3/2,89/64$.  The double-progression reduction bounds
every profile by ones with total at most $105/64$, and its displayed
profile is feasible.  Thus the reductions prove all five upper bounds,
while the listed feasible profiles attain them.

By Lemma~\ref{lem:cover}, each value is also the maximum of
$L_F(g,1_F)$ over nonzero nonnegative $g$, and the supremum when $g$
is required to be positive at every point.  This sharpness concerns the
universal translated covering problem.  The max-convolution ratio in
\eqref{eq:ratio} can be smaller than its covering bound.
